\begin{problem}[第35届全国中学生物理竞赛复赛][热辐射与热传导]{93}
六、如图，太空中有一由同心的内球和球壳构成的实验装置，内球和球壳内表面之间为真空。内球半径为 $r=0.200\ m$ ，温度保持恒定，比辐射率为 $e=0.800$；球壳的导热系数为 $\kappa=1.00\times10^{-2}\mathrm{J}\cdot\mathrm{m}^{-1}\cdot\mathrm{s}^{-1}\cdot\mathrm{K}^{-1}$ ，内、外半径分别为 $R_{1}=0.900\ m$ 、 $R_{2}=1.00\ m$ ，各表面的热辐射都是各向同性的，而外表面可视为黑体；该实验装置已处于热稳定状态，此时球壳内表面比辐射率为 $E=0.800$ 。斯特藩常量为 $\sigma=5.67\times10^{-8}W\cdot m^{-2}\cdot K^{-4}$ ，宇宙微波背景辐射温度为 $T=2.73K$ 。若单位时间内由球壳内表面传递到球壳外表面的热量为 $Q=44.0W$，求
\begin{subproblem}
球壳外表面温度 $T_{2}$ ；
\end{subproblem}
\begin{subproblem}
球壳内表面温度 $T_{1}$ ；
\end{subproblem}
\begin{subproblem}
内球温度 $T_{0}$ 。
\end{subproblem}
已知: 物体表面单位面积上的辐射功率与同温度下的黑体在该表面单位面积上的辐射功率之比称为比辐射率。当辐射照射到物体表面时，物体表面单位面积吸收的辐射功率与照射到物体单位面积上的辐射功率之比称为吸收比。物体在某一温度下的吸收比等于其在同一温度下的比辐射率。当物体内某处在 $z$ 方向（热流方向）每单位距离温度的增量为 $\frac{dT}{dz}$ 时，物体内该处单位时间在 $z$ 方向每单位面积流过的热量为 $-\kappa\frac{dT}{dz}$ ，此即傅里叶热传导定律。
\end{problem}

